\section*{第 \ref{ch:g8:solids} 章 --- 棱锥与圆锥}

\begin{solution}{exo:g8:solids:1}
四面体：$4$ 个面，$6$ 条棱，$4$ 个顶点
（$4 + 4 = 6 + 2$）。六棱锥：$7$ 个面，$12$ 条棱，$7$
个顶点（$7 + 7 = 12 + 2$）。
\end{solution}

\begin{solution}{exo:g8:solids:2}
$V = \frac13 \times 36 \times 10 = 120$ cm$^3$。

底面为长方形：$B = 24$ cm$^2$，于是
$V = \frac13 \times 24 \times 7.5 = 60$ cm$^3$。
\end{solution}

\begin{solution}{exo:g8:solids:3}
$V = \frac13 \pi \times 25 \times 9 = 75\pi \approx 236$ cm$^3$。
\end{solution}

\begin{solution}{exo:g8:solids:4}
圆柱：$\pi \times 16 \times 12 = 192\pi \approx 603$ cm$^3$。
圆锥：$\frac{192\pi}{3} = 64\pi \approx 201$ cm$^3$。两者之比：圆锥
的体积是圆柱的三分之一。
\end{solution}

\begin{solution}{exo:g8:solids:5}
$56 = \frac13 \times B \times 8$，于是 $B = \frac{56 \times 3}{8} = 21$
cm$^2$。
\end{solution}

\begin{solution}{exo:g8:solids:6}
侧面展开后是一个扇形，其半径为
\emph{斜高} $s$（从锥顶到底面圆周的距离 ---
该线段位于侧面上），而不是高 $h$。
\end{solution}

\begin{solution}{exo:g8:solids:7}
$V = \frac13 \times 35^2 \times 22 = \frac13 \times 1225 \times 22
= \frac{26\,950}{3} \approx 8\,983$ m$^3$ --- 约 $9\,000$ 立方
米。
\end{solution}

\begin{solution}{exo:g8:solids:8}
高：$h = \sqrt{13^2 - 5^2} = \sqrt{169 - 25} = \sqrt{144} = 12$
cm。体积：$V = \frac13 \pi \times 25 \times 12 = 100\pi$ cm$^3$。
\end{solution}

\begin{solution}{exo:g8:solids:9}
半径相同，因此圆锥的体积
是圆柱体积的三分之一（\emph{高度相同时}）：液面在圆柱中达到
$\frac{9}{3} = 3$ cm。
\end{solution}

\begin{solution}{exo:g8:solids:10}
沙子 $=$ 一个圆锥：$V = \frac13 \pi \times 4 \times 4.5 = 6\pi
\approx 18.8$ cm$^3$。沙漏装有两个这样的圆锥，因此
沙子恰好占据内部总体积的一半。
\end{solution}

\begin{solution}{exo:g8:solids:11}
\emph{1.} 是的：公式 $V = \frac13 Bh$ 对顶点的任意位置都成立，
只要 $h$ 是到底面
所在平面的垂直距离。$V = \frac13 \times 36 \times 8 = 96$ cm$^3$。

\emph{2.} 底面对角线：$\sqrt{6^2 + 6^2} = 6\sqrt2$ cm。最长的
侧棱是直角三角形的斜边，该直角三角形的两条直角边为
$6\sqrt2$（位于底面所在平面内）和 $8$（竖直方向）：
\[
\sqrt{(6\sqrt2)^2 + 8^2} = \sqrt{72 + 64} = \sqrt{136} \approx 11.7
\text{ cm}.
\]
\end{solution}

\begin{solution}{pb:g8:solids:1}
\textbf{1.} $G$ 属于上底面 $EFGH$，并属于两个侧面
$BCGF$ 和 $DCGH$。\emph{不}
含 $G$ 的三个面是下底面 $ABCD$、前侧面 $ABFE$ 和
左侧面 $ADHE$（它们都共用顶点 $A$，即与
$G$ 相对的顶点）。三个棱锥分别为 $ABCDG$、$ABFEG$ 和
$ADHEG$：每个都以正方体的一个正方形面为底，
以远端顶点 $G$ 为锥顶。每个都是四棱锥：$5$ 个面，
$8$ 条棱，$5$ 个顶点（\cref{ex:g8:solids:count}）。

\textbf{2.} 体对角线 $[AG]$ 连接的两个顶点，
不属于三个底面中的任何一个。绕直线
$(AG)$ 旋转三分之一周，使正方体与自身重合。
（离开 $A$ 的三条棱 --- 即指向 $B$、$D$、$E$ 的棱 --- 依次
轮换，指向 $G$ 的三条棱也同样轮换）。它把
底面 $ABCD$ 映到 $ADHE$，$ADHE$ 映到 $ABFE$，如此循环：
锥顶固定在 $G$，底面轮换。因此每个棱锥都被
映到下一个棱锥上，三者彼此全等。

\textbf{3.} 三个全等的部分填满正方体，
它们平分其体积：
\[
V_{\text{棱锥}} = \frac{a^3}{3} .
\]

\textbf{4.} 棱锥 $ABCDG$ 的底面为 $ABCD$，面积为 $a^2$，
其锥顶 $G$ 位于顶点 $C$ 的正上方，高为
$CG = a$（正方体的一条竖直棱）。与
\cref{exo:g8:solids:11}一样，公式取垂直高度
$h = a$：
\[
V = \frac13 \times a^2 \times a = \frac{a^3}{3} ,
\]
与第 3 小题完全一致 --- 分解与
公式互相印证。

\textbf{5.} 取 $a = 6$：
每个棱锥的体积为 $V = \frac{6^3}{3} = \frac{216}{3} = 72$ cm$^3$。
倒水实验就是这一分解的液体版本：
以 $ABCD$ 为底、高为 $a$ 的棱柱即正方体本身，
它恰好容纳三份棱锥的水量 --- 这里的三份
甚至在几何上拼成了正方体。

\textbf{6.} 将中心 $O$ 与每个面的四个顶点相连，
由此得到六个棱锥，每个面一个：底面为正方形，顶点为 $O$。
中心与六个面的位置关系相同（立方体的任一旋转
或对称都使 $O$ 不动，并置换各个面），
因此六个棱锥全等。每个棱锥的高为
$O$ 到一个面的距离，即棱长的一半，$\frac{a}{2}$。

\textbf{7.} 六块全等的部分拼成正方体：
每块 $V = \frac{a^3}{6}$。

\textbf{8.} 由公式，底面积为 $a^2$，高为 $\frac a2$：
\[
V = \frac13 \times a^2 \times \frac{a}{2} = \frac{a^3}{6} .
\]
两个结果一致。

\textbf{9.} 两种分解所涉及的都是从正方体中切出的、比例十分特殊的棱锥；
细长的棱锥、歪斜的棱锥或圆锥都不能这样拼出
（没有三个圆锥能充满一个圆柱）。这正是
\cref{thm:g8:solids:volume}在本阶段
\emph{不加证明地承认}：
分解与倒水实验使 $\frac13$ 变得完全可信，
而严谨的一般证明 --- 积分 --- 在
高中分册中给出。

\textbf{10.} $V = \frac{6^3}{6} = 36$ cm$^3$；由公式，
$V = \frac13 \times 36 \times 3 = 36$ cm$^3$。两者一致。

\textbf{11.} 在三角形 $SAB$ 中，边 $[SA]$ 与 $[SB]$ 的中点
由一条与 $(AB)$ 平行的线段相连，
且其长度为 $\frac{AB}{2} = \frac{c}{2}$
（见\cref{thm:g8:midpoints:direct}）。每个侧面上的情形相同，
因此四个中点构成边长为 $\frac c2$ 的正方形，
各边与底面平行。过顶点的竖直线上，
截面恰好经过高的中点（再次用到中点定理，
在经过 $S$、底面中心与底面一个顶点的三角形中）；
即截面所在平面的高为 $\frac h2$。

\textbf{12.} 上方的小块是一个正四棱锥，
底面边长为 $\frac c2$、高为 $\frac h2$（其底面即为截面，
其顶点为 $S$）。其体积为
\[
\frac13 \times \left(\frac c2\right)^2 \times \frac h2
= \frac13 \times \frac{c^2}{4} \times \frac h2
= \frac{1}{8} \times \frac{c^2 h}{3}
= \frac{V}{8} :
\]
即原体积的八分之一 --- 而不是一半。

\textbf{13.} 剩余部分是棱台：
$V - \frac V8 = \frac{7V}{8}$。半高截面将体积按
$1 : 7$ 的比例分配 --- 下方看似其貌不扬的棱台
所占体积是上方尖顶部分的七倍。

\textbf{14.} 半高以下的部分占体积的 $\frac78$。取
$V = \frac13 \times 35^2 \times 22 = \frac{26\,950}{3} \approx
8\,983$ m$^3$（见\cref{exo:g8:solids:7}）：
\[
\frac78 \times \frac{26\,950}{3} \approx 7\,860 \approx
7\,900 \text{ m}^3
\]
精确到百位 --- “下半部分”行动
清洗的玻璃体积几乎是上半部分的八倍。

\textbf{15.} 将各尺寸加倍，底面积变为原来的
$2^2 = 4$ 倍，高变为原来的 $2$ 倍，因此体积变为原来的
$4 \times 2 = 8 = 2^3$ 倍。若按 $k$ 倍缩放，面积变为原来的 $k^2$
倍，另一长度变为原来的 $k$ 倍：体积按 $k^3$ 的因子增大。
第 12 题中的上方小块是整个棱锥的相似形，
相似比为 $k = \frac12$，因此其体积为
整体的 $\left(\frac12\right)^3 = \frac18$ --- 这便是
一句话的解释。
\end{solution}
